\chapter*{این کتابچه چیست}
\addcontentsline{toc}{chapter}{این کتابچه چیست}
\markboth{این کتابچه چیست}{این کتابچه چیست}

\textit{متسیس} یک کتاب تمام‌شده نیست. ثبتِ جاریِ تلاش یک نفر است برای
فهمیدن ریاضیات از ریشه‌هایش --- از منطق و مجموعه‌ها به بعد --- که
کتابچه به کتابچه، همان‌طور که فهم واقعاً ساخته می‌شود مستند می‌شود،
نه این‌که بعداً از موضع کسی که از قبل می‌دانسته خلاصه شده باشد.

نسخه‌ی فارسیِ هر کتابچه فعلاً فقط شامل همین بخش‌های ساختاری است؛
فصل‌های اصلی به‌مرور و هم‌زمان با نسخه‌ی انگلیسی نوشته و اضافه
می‌شوند.

\begin{notebox}[title=یادداشت]
جعبه‌هایی مثل این، نکته‌ای را که متن اصلی بدون توضیح مستقیم به آن
تکیه می‌کند مشخص می‌کنند.
\end{notebox}

\begin{tipbox}[title=نکته]
این قالب یک پیشنهاد عملی را نشان می‌دهد --- روشی دیگر برای خواندن یک
بخش، یا اشاره‌ای به چیزی که ارزش پیگیری بیشتر دارد.
\end{tipbox}
